% ---------------------------------------------------------------------
% 2.4 Interfaces cerebro-computadora
% ---------------------------------------------------------------------
\section[INTERFACES CEREBRO-COMPUTADORA]{Interfaces cerebro-computadora}
\label{sec:bci}

\subsection{Definición y componentes}

Las interfaces cerebro-computadora (BCI) permiten entablar comunicación directa entre la
actividad cerebral de un individuo y un dispositivo externo, sin depender de ninguna otra señal o
movimiento muscular. De acuerdo con Wolpaw et al., una BCI es un sistema de comunicación en el que
los mensajes o comandos que el usuario envía al mundo exterior no pasan por las vías de salida
normales del cerebro, es decir, por los nervios periféricos y los músculos \cite{wolpaw2002}, y
como se menciona en el trabajo de Felton et al., estos sistemas traducen las señales cerebrales
registradas en acciones en tiempo real, lo cual abre la posibilidad de restaurar funciones
comunicativas o motoras en personas con discapacidad severa \cite{felton2007}.

Independientemente de la señal que utilice, una BCI se compone de una serie de etapas que se
ejecutan de forma cíclica, como se ilustra en la Figura~\ref{fig:bci}. Primero se adquiere y digitaliza la
señal cerebral. Después se preprocesa para reducir el ruido y los artefactos. Luego se extraen las
características que mejor representan la intención del usuario, y un clasificador las convierte en
un comando para el dispositivo. Finalmente, el resultado se muestra al usuario como
retroalimentación para que pueda ajustar su estrategia \cite{wolpaw2012,nicolasalonso2012}.

\begin{figure}[htbp]
\centering
\begin{tikzpicture}[
  node distance=0.38cm,
  etapa/.style={draw, rounded corners=2pt, fill=gray!8, align=center, minimum height=1.1cm, text width=2.15cm, font=\footnotesize},
  flecha/.style={-{Stealth[length=2mm]}, thick}]
  \node[etapa, fill=red!8] (usr) {Usuario\\(actividad cerebral)};
  \node[etapa, right=of usr] (adq) {Adquisición\\(EEG)};
  \node[etapa, right=of adq] (pre) {Preprocesa-\\miento};
  \node[etapa, right=of pre] (car) {Extracción de\\características};
  \node[etapa, right=of car] (cla) {Clasificación};
  \node[etapa, below=1.1cm of cla, fill=blue!6] (app) {Aplicación\\(deletreador)};
  \draw[flecha] (usr) -- (adq);
  \draw[flecha] (adq) -- (pre);
  \draw[flecha] (pre) -- (car);
  \draw[flecha] (car) -- (cla);
  \draw[flecha] (cla) -- node[right, font=\footnotesize] {comando} (app);
  \draw[flecha, dashed] (app.west) -| node[pos=0.25, below, font=\footnotesize] {retroalimentación visual} (usr.south);
  \begin{scope}[on background layer]
    \node[draw=gray, dashed, rounded corners, fit=(pre)(car)(cla), inner sep=6pt,
          label={[font=\footnotesize]above:Procesamiento de la señal}] {};
  \end{scope}
\end{tikzpicture}
\caption{Diagrama general de una interfaz cerebro-computadora. Elaboración propia con base en
\cite{wolpaw2002,wolpaw2012}.}
\label{fig:bci}
\end{figure}

\subsection{Clasificación de las BCI}

Las BCI pueden clasificarse bajo distintos criterios. Según la forma en que se registra la señal,
se dividen en invasivas, que requieren colocar electrodos dentro o sobre el cerebro, y no
invasivas, que registran la actividad desde el exterior \cite{nicolasalonso2012}. Según el origen
de la actividad que aprovechan, se distinguen las exógenas, que dependen de una estimulación
externa como la iluminación de una matriz de caracteres, y las endógenas, en las que el usuario
modula voluntariamente su actividad cerebral sin estímulo alguno, por ejemplo al imaginar el
movimiento de una mano \cite{nicolasalonso2012}. Asimismo, una BCI es síncrona cuando el sistema
determina los intervalos en los que el usuario puede emitir un comando, y asíncrona cuando el
usuario puede hacerlo en cualquier momento, lo que exige detectar también los periodos en los que
no desea comunicarse \cite{wolpaw2012}. Bajo estos criterios, el P300 Speller es una BCI no
invasiva, exógena y síncrona.

\subsection{Paradigmas de control}

Las BCI basadas en EEG utilizan principalmente cuatro tipos de señales. Las basadas en ritmos sensoriomotores aprovechan la
disminución de los ritmos mu y beta sobre la corteza motora al imaginar un movimiento. Las de
potenciales evocados visuales de estado estable (SSVEP) usan la respuesta de la corteza visual a
estímulos que parpadean a una frecuencia fija. Las de potenciales corticales lentos requieren que
el usuario aprenda a modificar voluntariamente la línea base de su EEG. Por último, las basadas en
la P300 aprovechan la respuesta del paradigma del evento raro descrita en la sección anterior
\cite{nicolasalonso2012,wolpaw2002}. Una de las principales ventajas de este último enfoque es que
la P300 es una respuesta natural ante los estímulos relevantes, por lo que prácticamente no
requiere entrenamiento previo del usuario \cite{farwell1988,contreras2021}.
